\section{Билет 1. Задача ЛП в канонической форме}

\ticketq{Задача ЛП в канонической форме. Основные понятия. Лемма о числе
опорных векторов. Теорема о совпадении крайней точки и опорного вектора.
Теорема об оптимальном решении задачи ЛП. Метод перебора крайних точек.}

\paragraph{Постановка.}
\[
  c\T[N]x[N] \to \min , \qquad A[M,N]x[N] = b[M] , \qquad x[N] \ge 0 ,
\]
$\abs{M} = m$, $\abs{N} = n$, $m < n$, $\rank A = m$;
$S = \{\, x[N] \mid A[M,N]x[N] = b[M],\ x[N] \ge 0 \,\}$ — многогранное
допустимое множество.

\paragraph{Основные понятия.}
$x_0 \in S$ — \textbf{крайняя точка}, если
$x_0 \ne (1-\lambda)x_1 + \lambda x_2$ ни при каких $x_1 \ne x_2 \in S$,
$0 < \lambda < 1$.

$\Nzp = \{j : x_0[j] > 0\}$, $\Nzz = \{j : x_0[j] = 0\}$
$\Rightarrow$ $\sum_{j \in \Nzp} x_0[j]\,a_j[M] = b[M]$ ($a_j$ — столбцы $A$).

$x_0$ — \textbf{опорный}, если $\{a_j[M]\}_{j \in \Nzp}$ линейно независимы;
они в $\R^m$ $\Rightarrow$ $\abs{\Nzp} \le m$. Невырожденный:
$\abs{\Nzp} = m$. Вырожденный: $\abs{\Nzp} < m$, тогда $\Nzp \subset N_0$,
$\abs{N_0} = m$, $\det A[M,N_0] \ne 0$, $x_0[N_0 \setminus \Nzp] = 0$.
\textbf{Базис} — столбцы $a_j$, $j \in N_0$.

\paragraph{Лемма (о числе опорных векторов).}
\textit{Опорных векторов не больше $\binom{n}{m}$.}

\textsc{Док-во.} Опорный вектор $\mapsto$ базис $=$ набор $m$ столбцов из $n$;
таких наборов $\binom{n}{m}$. Для набора $N_k$ с $\det A[M,N_k] \ne 0$ решение
системы $A[M,N_k]x[N_k] = b[M]$ единственно и дополняется нулями
$\Rightarrow$ разных опорных векторов не больше числа наборов.
$\blacksquare$

\paragraph{Теорема (крайняя точка $=$ опорный вектор).}
\textit{$x[N] \in S$ опорный $\iff$ $x[N]$ — крайняя точка $S$.}
Далее $\Np = \{j : x[j] > 0\}$.

\textsc{Док-во. 1) Не опорный $\Rightarrow$ не крайняя.}
\begin{gather*}
  \exists\, u[\Np] \ne 0 :\ A[M,\Np]u[\Np] = 0 , \qquad
  u[N \setminus \Np] := 0 \\
  \Rightarrow\quad A[M,N]u[N] = 0
  \quad\Rightarrow\quad A(x \pm \theta u) = b .
\end{gather*}
Неотрицательность держится при $\theta \le \theta_0$:
\[
  \theta_0 = \min\{\theta_1, \theta_2\} , \quad
  \theta_1 = \min_{i \in \Np,\, u[i]>0} \frac{x[i]}{u[i]} , \quad
  \theta_2 = \min_{i \in \Np,\, u[i]<0} \left(-\frac{x[i]}{u[i]}\right) .
\]
$x_{1,2} = x \pm \theta_0 u \in S$, $x_1 \ne x_2$ (ибо $u \ne 0$),
$x = \tfrac12(x_1+x_2)$ $\Rightarrow$ $x$ не крайняя.

\textsc{2) Не крайняя $\Rightarrow$ не опорный.}
$x = (1-\lambda)x_1 + \lambda x_2$. Для $i \in \Nz$:
\[
  0 = \underbrace{(1-\lambda)x_1[i]}_{\ge\,0}
    + \underbrace{\lambda x_2[i]}_{\ge\,0}
  \ \Rightarrow\ x_1[i] = x_2[i] = 0
\]
$\Rightarrow$ $x_1, x_2$ различаются лишь на $\Np$. Из $Ax_1 = Ax_2 = b$:
\[
  A[M,\Np]\big(x_1[\Np] - x_2[\Np]\big) = 0 ,
  \quad x_1[\Np] - x_2[\Np] \ne 0
  \ \Rightarrow\ \text{столбцы зависимы.} \qquad \blacksquare
\]

\paragraph{Теорема (об оптимальном решении).}
\textit{Если задача ЛП разрешима, то хотя бы одно оптимальное решение —
опорный вектор.}

\textsc{Док-во.} Берём оптимальное $\bar{x}$ с минимальным
$\abs{\overline{N}{}^+}$, где $\overline{N}{}^+ = \{j : \bar{x}[j] > 0\}$.
Пусть $\bar{x}$ не опорный:
\[
  \exists\, u[\overline{N}{}^+] \ne 0 :\ A[M,\overline{N}{}^+]u = 0 ,
  \quad u[N \setminus \overline{N}{}^+] := 0
  \ \Rightarrow\ A[M,N]u[N] = 0 .
\]
Тогда $x_\theta := \bar{x} + \theta u$ даёт $Ax_\theta = b$ и $x_\theta \ge 0$
при $\abs{\theta} \le \theta_0$, то есть $x_\theta \in S$, причём
\[
  c\T x_\theta = c\T\bar{x} + \theta\, c\T u
  \ \Rightarrow\
  \big[\, c\T u \ne 0 \ \Rightarrow\ \exists\,\theta :
    c\T x_\theta < c\T\bar{x} \,\big]
  \ —\ \text{противоречие.}
\]
Значит $c\T u = 0$ и все $x_\theta$ оптимальны. При $\theta = \theta_1$
(либо $-\theta_2$) одна положительная компонента обнуляется, новых не
возникает ($u = 0$ вне $\overline{N}{}^+$)
$\Rightarrow$ оптимум с меньшим числом положительных компонент —
противоречие с выбором $\bar{x}$. $\blacksquare$

\textbf{Следствие.} Оптимум достигается в крайней точке; их конечное число.

\paragraph{Метод перебора крайних точек.}
\begin{steps}
  \item Взять $N_k \subset N$, $\abs{N_k} = m$.
  \item $\det A[M,N_k] = 0$ — отбросить.
  \item Решить $A[M,N_k]x[N_k] = b[M]$, положить $x[N \setminus N_k] = 0$.
  \item $x[N_k] \ge 0$ — крайняя точка, считаем $F = c\T[N]x[N]$; иначе
        отбросить.
  \item Перебрать все наборы, взять минимальное $F$.
\end{steps}

Конечен по лемме; нет допустимых точек $\Rightarrow$ $S = \varnothing$.
Непригоден практически: $\binom{50}{25} \approx 1{,}3\cdot10^{14}$, и
неограниченность не распознаётся $\Rightarrow$ симплекс-метод.

\seefull{гл.~1, \S1--\S6 и \S8}
